\documentclass[11pt]{article}
\usepackage[utf8]{inputenc}
\usepackage[spanish]{babel}
\usepackage[margin=2.5cm]{geometry}
\usepackage{titlesec}
\usepackage{enumitem}
\usepackage{fancyhdr}
\usepackage{xcolor}
\usepackage{tabularx}

\definecolor{ucrblue}{RGB}{20,60,90}

\titleformat{\section}{\normalfont\Large\bfseries\color{ucrblue}}{\thesection}{1em}{}
\titleformat{\subsection}{\normalfont\large\bfseries}{\thesubsection}{1em}{}

\pagestyle{fancy}
\fancyhf{}
\renewcommand{\headrulewidth}{0pt}
\cfoot{\thepage}

\newenvironment{sesion}[2]{%
  \subsection*{Fecha: #1 \hfill Horas invertidas: #2}
  \vspace{-0.3cm}
  \hrule
  \vspace{0.3cm}
}{\vspace{0.5cm}}

\newcommand{\campo}[1]{\textbf{#1}}

\begin{document}

% ---------- PORTADA ----------
\begin{titlepage}
\centering
\vspace*{1.5cm}
{\large Universidad de Costa Rica\par}
{\large Escuela de Ingeniería Eléctrica\par}
\vspace{1cm}
{\Large Proyecto Eléctrico -- IE-0499\par}
\vspace{2cm}
{\huge \bfseries Bitácora\par}
\vspace{1cm}
\rule{0.7\textwidth}{0.4pt}\\[0.4cm]
{\Large\bfseries Detección en tiempo real de deepfakes de audio\\ (voces sintéticas) aplicada al español costarricense\par}
\rule{0.7\textwidth}{0.4pt}
\vspace{3cm}

{\large Gabriel Esquivel Núñez\par}
{\large Carné: C22799\par}

\vfill

{\large II Ciclo -- 2026\par}
\vspace{0.5cm}
\end{titlepage}

% ---------- DATOS GENERALES DEL PROYECTO ----------
\section*{Datos generales del proyecto}

\campo{Título:} Detección en tiempo real de deepfakes de audio (voces sintéticas) aplicada al español costarricense.

\vspace{0.3cm}
\campo{Objetivo general:} Implementar y evaluar un sistema de detección de voces sintéticas (deepfakes de audio) capaz de operar sobre audio continuo o en flujo, mediante la adaptación de técnicas existentes de anti-spoofing, con aplicación a voces en español costarricense.

\vspace{0.3cm}
\campo{Objetivos específicos:}
\begin{enumerate}[label=\arabic*.]
\item Revisar el estado del arte en detección de deepfakes de audio y anti-spoofing de voz, con énfasis en técnicas adaptables al procesamiento por ventanas y a escenarios de tiempo real.
\item Seleccionar e implementar al menos dos enfoques de detección complementarios (uno basado en características acústicas y otro en modelos o representaciones preentrenadas), integrados a la estructura del repositorio del laboratorio.
\item Adaptar el detector para procesar audio en ventanas temporales sucesivas y estudiar el efecto de la duración de ventana y del mecanismo de decisión temporal sobre precisión, latencia y tiempo de decisión confiable.
\item Evaluar el desempeño y la robustez del sistema sobre voces naturales y sintéticas, y construir un demostrador funcional de detección en reproducción o captura continua de audio.
\end{enumerate}

\vspace{0.3cm}
\campo{Entregables:}
\begin{enumerate}[label=\arabic*.]
\item Informe de revisión y diseño del sistema de detección en tiempo real -- 4 de setiembre de 2026.
\item Detector base y procesamiento por ventanas -- 25 de setiembre de 2026.
\item Sistema de detección en tiempo real e informe de evaluación -- 24 de octubre de 2026.
\item Informe final y demostrador funcional -- 27 de noviembre de 2026.
\end{enumerate}

\vspace{0.3cm}
\campo{Proyecto base del laboratorio:} el trabajo parte de la infraestructura del proyecto \textit{voice-similarity-eval} del laboratorio (LAPA-CR), un toolkit reproducible para evaluar similitud de locutor y calidad perceptual en sistemas TTS y de conversión de voz. El presente proyecto complementa dicho toolkit con un módulo de detección de audio sintético (anti-spoofing) orientado a operación en tiempo real.

\newpage

% ---------- SEMANA 1 ----------
\section{Semana 1 (10/08/2026 -- 16/08/2026)}

\begin{sesion}{11/08/2026}{2 horas}
\campo{Objetivos.}
\begin{itemize}
\item Confirmar el alcance del proyecto y los entregables definidos en la propuesta.
\item Identificar los conceptos clave que era necesario investigar antes de iniciar la implementación.
\end{itemize}
\campo{Actividades y resultados.}
\begin{itemize}
\item Se revisó a fondo el enunciado de la propuesta de proyecto eléctrico, identificando los cuatro entregables y sus fechas límite.
\item Se listaron los conceptos a investigar: anti-spoofing de voz, retos ASVspoof, características acústicas (CQCC/LFCC), modelos end-to-end (RawNet2, AASIST) y representaciones autosupervisadas (wav2vec2, WavLM).
\end{itemize}
\campo{Dificultades encontradas.} Ninguna relevante; la sesión fue principalmente de planificación.
\end{sesion}

\begin{sesion}{13/08/2026}{3 horas}
\campo{Objetivos.}
\begin{itemize}
\item Comprender el problema de anti-spoofing en verificación de locutor y su relación con los sistemas de detección de deepfakes de audio.
\item Estudiar los enfoques clásicos basados en características acústicas.
\end{itemize}
\campo{Actividades y resultados.}
\begin{itemize}
\item Se investigó el marco general de los retos ASVspoof (incluyendo la edición más reciente, ASVspoof5) como referencia estándar del área.
\item Se estudiaron los coeficientes CQCC y LFCC como características acústicas interpretables, y su combinación típica con clasificadores como SVM o regresión logística.
\end{itemize}
\campo{Dificultades encontradas.} La bibliografía asume conocimientos previos de procesamiento de señales de voz (transformada Q constante, análisis cepstral), por lo que se requirió tiempo adicional de repaso.
\end{sesion}

\begin{sesion}{15/08/2026}{3 horas}
\campo{Objetivos.}
\begin{itemize}
\item Estudiar los modelos de aprendizaje profundo de referencia en anti-spoofing (RawNet2, AASIST) y las representaciones autosupervisadas (wav2vec2, WavLM).
\item Revisar la estructura y el alcance del repositorio base del laboratorio.
\end{itemize}
\campo{Actividades y resultados.}
\begin{itemize}
\item Se investigaron las arquitecturas de RawNet2 y AASIST, así como la disponibilidad de pesos preentrenados sobre ASVspoof para uso por transferencia.
\item Se revisó el repositorio \textit{voice-similarity-eval} (LAPA-CR), utilizado por el laboratorio para evaluar similitud de locutor y calidad perceptual en voces sintéticas, con el fin de entender su estructura modular (\texttt{features/}, \texttt{metrics/}, \texttt{perceptual\_tests/}, \texttt{robustness/}) y cómo se integraría un módulo de detección de audio sintético.
\end{itemize}
\campo{Dificultades encontradas.} El repositorio base está orientado a evaluación de similitud, no a detección de spoofing, por lo que fue necesario analizar cómo adaptar su estructura modular a un módulo nuevo de anti-spoofing en tiempo real.
\end{sesion}

% ---------- SEMANA 2 ----------
\section{Semana 2 (17/08/2026 -- 23/08/2026)}

\begin{sesion}{18/08/2026}{2 horas}
\campo{Objetivos.}
\begin{itemize}
\item Investigar las consideraciones necesarias para adaptar los detectores a un régimen de ventanas cortas de audio.
\end{itemize}
\campo{Actividades y resultados.}
\begin{itemize}
\item Se estudiaron las estrategias de agregación temporal de decisiones (decisión independiente por ventana, promedio de probabilidades, voto mayoritario, umbral sostenido).
\item Se identificó el compromiso entre duración de ventana, latencia y confiabilidad de la detección como uno de los principales retos técnicos del proyecto.
\end{itemize}
\campo{Dificultades encontradas.} Ninguna relevante.
\end{sesion}

\begin{sesion}{20/08/2026}{3 horas}
\campo{Objetivos.}
\begin{itemize}
\item Definir la estructura del informe de revisión y diseño del sistema (primer entregable).
\end{itemize}
\campo{Actividades y resultados.}
\begin{itemize}
\item Se definió la estructura del documento: introducción, estado del arte, selección justificada de enfoques, protocolo de evaluación, arquitectura preliminar y conclusiones.
\item Se creó el esqueleto del documento en LaTeX, formato IEEEtran, con la bibliografía inicial de referencia (ASVspoof5, RawNet2, AASIST, wav2vec2, WavLM, CQCC).
\end{itemize}
\campo{Dificultades encontradas.} Ninguna relevante.
\end{sesion}

\begin{sesion}{22/08/2026}{2 horas}
\campo{Objetivos.}
\begin{itemize}
\item Crear el repositorio de trabajo para el proyecto y dejarlo listo para comenzar la fase de implementación.
\end{itemize}
\campo{Actividades y resultados.}
\begin{itemize}
\item Se creó el repositorio de trabajo del proyecto, tomando como referencia la estructura del repositorio base del laboratorio (\textit{voice-similarity-eval}).
\item Se organizó una carpeta inicial para el módulo de detección de audio sintético, coherente con la convención modular del repositorio base (\texttt{features/}, \texttt{metrics/}, etc.).
\end{itemize}
\campo{Dificultades encontradas.} Ninguna relevante.
\end{sesion}

% ---------- SEMANA 3 ----------
\section{Semana 3 (24/08/2026 -- 30/08/2026)}

\begin{sesion}{25/08/2026}{3 horas}
\campo{Objetivos.}
\begin{itemize}
\item Redactar la sección de estado del arte del informe de revisión.
\end{itemize}
\campo{Actividades y resultados.}
\begin{itemize}
\item Se redactó la sección de estado del arte, cubriendo enfoques basados en características acústicas (CQCC/LFCC), modelos end-to-end (RawNet2, AASIST) y representaciones autosupervisadas (wav2vec2, WavLM).
\item Se incluyó una subsección específica sobre las consideraciones necesarias para adaptar estos enfoques a un escenario de operación en tiempo real por ventanas.
\end{itemize}
\campo{Dificultades encontradas.} Ninguna relevante.
\end{sesion}

\begin{sesion}{27/08/2026}{2 horas}
\campo{Objetivos.}
\begin{itemize}
\item Justificar la selección de los dos enfoques de detección a implementar.
\item Definir el protocolo de evaluación del proyecto.
\end{itemize}
\campo{Actividades y resultados.}
\begin{itemize}
\item Se seleccionó y justificó el enfoque interpretable (LFCC + SVM/regresión logística) y el enfoque preentrenado (AASIST por transferencia) como las dos técnicas complementarias a implementar.
\item Se definió el protocolo de evaluación: reutilización del conjunto de voces del laboratorio, particiones sin fuga de información entre locutores/sistemas TTS, longitudes de ventana a evaluar (0.5, 1, 2, 3 y 5 s) y métricas de clasificación y de tiempo real (EER, F1, latencia, RTF).
\end{itemize}
\campo{Dificultades encontradas.} Ninguna relevante.
\end{sesion}

\begin{sesion}{29/08/2026}{3 horas}
\campo{Objetivos.}
\begin{itemize}
\item Diseñar la arquitectura preliminar del sistema de detección continua.
\item Finalizar y compilar el informe de revisión y diseño.
\end{itemize}
\campo{Actividades y resultados.}
\begin{itemize}
\item Se diseñó el diagrama de bloques de la arquitectura preliminar: entrada de audio, segmentación en ventanas, procesamiento en paralelo por los dos detectores, acumulación temporal de decisiones y salida real/falsa con nivel de confianza.
\item Se completó la redacción de las conclusiones preliminares del informe y se compiló el documento final en PDF, quedando listo para su revisión antes de la entrega del 4 de setiembre.
\end{itemize}
\campo{Dificultades encontradas.} Ninguna relevante.
\end{sesion}

\end{document}
